\documentclass[11pt]{article}
\usepackage[margin=1in]{geometry}
\usepackage{amsmath,amssymb,amsthm,booktabs,enumitem,microtype}
\usepackage[colorlinks=true,linkcolor=blue,urlcolor=blue,citecolor=blue]{hyperref}

\setlist{leftmargin=*,itemsep=3pt}
\newcommand{\fatal}{\textbf{Fatal}}
\newcommand{\major}{\textbf{Major}}
\newcommand{\moderate}{\textbf{Moderate}}
\newcommand{\minor}{\textbf{Minor}}

\title{Adversarial Mathematical Assessment of \texttt{sumdist.tex}}
\author{}
\date{12 July 2026}

\begin{document}
\maketitle

\begin{abstract}
This report audits the manuscript \texttt{sumdist.tex}, which claims the
order-sharp asymptotic
\[
 \frac43N^2-\sum_{i\ne j}|x_i-x_j|\asymp\sqrt N
\]
for the spherical point set constructed by Beltr\'an, Etayo, Marzo and
Ortega-Cerd\`a (BEMOC).  The exact spectral and ring decompositions are sound,
as is the calculation of the within-ring contribution.  The claimed upper
bound is not, however, completely proved: the cross-ring Fourier estimate is
not uniform in the coalescing-branch-point regime, and the necklace section
contains invalid summations and an $N^2$ scaling error.  The stronger claims
about the necklace constant rely on unproved numerical shell estimates.  We
state precisely which results survive and describe a viable route to repair
the main $O(\sqrt N)$ estimate.
\end{abstract}

\section{Summary verdict}

Write
\[
 \Sigma_N=\sum_{i\ne j}|x_i-x_j|,
 \qquad D_N=\frac43N^2-\Sigma_N.
\]
The manuscript does \emph{not} currently prove $D_N\asymp\sqrt N$.

\begin{center}
\begin{tabular}{p{0.30\textwidth}p{0.16\textwidth}p{0.44\textwidth}}
\toprule
Component & Verdict & Reason\\
\midrule
Legendre expansion and sign & Correct & Coefficients and addition-theorem argument check.\\
Exact $A+B+C$ decomposition & Correct & Algebraic rearrangement and spectral signs check.\\
Within-ring term $B$ & Correct & Exact cotangent identity and asymptotic are valid.\\
Cross-ring bound $|C|\lesssim\sqrt N$ & Major gap & Required Fourier bound is not proved uniformly as branch points approach the real axis.\\
Necklace bound $A\lesssim\sqrt N$ & Fatal gap & Residual summation and far-band scaling are invalid.\\
Necklace lower bound and limit & Unproved & ``Few percent'' shell estimate is asserted numerically, without uniform proof.\\
Universal lower bound $D_N\gtrsim\sqrt N$ & Correct & Stolarsky--Beck argument applies.\\
\bottomrule
\end{tabular}
\end{center}

Thus the manuscript rigorously supplies the exact decomposition, the
universal lower bound, and the asymptotic of $B$, but not the advertised
matching upper bound.

\section{Verified components}

\subsection{Spectral identity}

The Legendre expansion in lines 309--348 is correct:
\[
 \sqrt{2-2u}
 =\frac43-\sum_{n\ge1}\frac4{(2n-1)(2n+3)}P_n(u).
\]
The coefficients are absolutely summable, so termwise summation over the
ordered pairs is justified.  The addition theorem gives
\[
 S_n:=\sum_{i,j}P_n(\langle x_i,x_j\rangle)\ge0,
\]
and hence
\[
 \Sigma_N=\frac43N^2-\sum_{n\ge1}\frac4{(2n-1)(2n+3)}S_n,
 \qquad D_N\ge0.
\]
The ordered-pair convention is harmless because the diagonal distances
vanish.

\subsection{Ring decomposition and signs}

The identity in lines 403--408,
\[
 D_N=A+B+C,
\]
is an exact algebraic decomposition obtained by inserting the continuous
angular measure on every occupied parallel.  Its spectral interpretation is
also correct:
\[
 A\ge0,\qquad B+C\ge0.
\]
The cross-ring term $C$ is signed; the manuscript does not incorrectly claim
that $B$ and $C$ are separately nonnegative.

\subsection{Within-ring contribution}

For a ring of radius $\rho$ carrying $w$ equally spaced points, the exact
ordered energy is
\[
 2\rho w\cot\frac{\pi}{2w}.
\]
Its deficit relative to the continuous angular measure is
\[
 \rho\left(\frac{4w^2}{\pi}-2w\cot\frac{\pi}{2w}\right)
 =\frac{\pi\rho}{3}+O\left(\frac{\rho}{w^2}\right).
\]
Consequently the calculation
\[
 B=\frac{\pi}{3}\sum_p\rho_p+o(1)
\]
is consistent with the BEMOC ring geometry.  Along $N=4m^2$,
\[
 \frac1{\sqrt N}\sum_p\rho_p\longrightarrow
 \frac{2(2\sqrt2-1)}3,
\]
so
\[
 \frac{B}{\sqrt N}\longrightarrow
 c_B:=\frac{2\pi(2\sqrt2-1)}9
 =1.276482938\ldots.
\]
The decimals $1.2761$ and $1.2766$ in the manuscript should be corrected.

\subsection{Universal lower bound}

With normalized surface area and ordinary Lebesgue measure $dt$ on
$[-1,1]$, Stolarsky's identity reads
\[
 \frac43-\frac{\Sigma_N}{N^2}
 =4\int_{-1}^1\int_{\mathbb S^2}
 \left|\frac1N\sum_{i=1}^N\mathbf1_{\{x_i\cdot z\ge t\}}
 -\sigma\{x:x\cdot z\ge t\}\right|^2d\sigma(z)\,dt.
\]
The coefficient is $4$ in this normalization.  Beck's spherical-cap
$L^2$ discrepancy bound implies that the double integral is at least
$cN^{-3/2}$ and therefore
\[
 D_N\ge c\sqrt N.
\]

\section{Cross-ring term}

\subsection{Correct Poisson identity}

For two rings of populations $q,r$, relative phase $\delta$, and
\[
 G(\theta)=\sqrt{a-b\cos\theta},
\]
the Fourier calculation in lines 508--541 is correct.  Only modes divisible
by $L=\operatorname{lcm}(q,r)$ survive, giving
\[
 qr I-\Sigma^{\rm cross}
 =-qr\sum_{\substack{m\ne0\\L\mid m}}widehat G_m e^{im\delta}.
\]

\subsection{Nonuniform branch-point estimate}

\major{} Lines 559--568 claim
\[
 |\widehat G_m|le C_0\sqrt{b\sinh\beta}\,
 m^{-3/2}e^{-m\beta}
\]
with one absolute constant.  The cited Watson/singularity-analysis arguments
give fixed-parameter asymptotics.  They do not by themselves give a bound
uniform when
\[
 \beta\to0,\qquad b\to0,
 \qquad m=L\to\infty,
 \qquad L\beta=O(1).
\]
This is precisely the adjacent-ring regime needed in the proof.

Further omissions are:
\begin{itemize}
\item Rings immediately north and south of the equator are not exponentially
suppressed, contrary to lines 580--582.
\item The text substitutes $L\asymp j'$ although only $L\gtrsim j'$ is known.
Monotonicity can make this repairable, but it must be stated.
\item The BEMOC rings comprise two interleaved families.  The conversion from
band index to one globally uniformly spaced ring index is not proved.
\item Uniform first-term domination of the aliasing series is not established
for cross-equatorial neighboring rings.
\end{itemize}

Thus the proof of $|C|\lesssim\sqrt N$ is incomplete.

\subsection{A viable repair}

There is a simpler parameter-uniform route.  If
$G^2=A-B\cos\theta$ with $A\ge B$, then
\[
 \|DG'\|_{\rm TV}\le C\sqrt B.
\]
Indeed, $G'$ is bounded by $C\sqrt B$, and the equation $G''=0$ has uniformly
boundedly many solutions.  Distributionally this remains true in the cusp
case $A=B$.  Two integrations by parts give
\[
 |\widehat G_n|\le\frac{C\sqrt B}{n^2}.
\]
The exact least-common-multiple reduction then yields the phase-uniform
pairwise error
\[
 |\operatorname{error}(q,r)|
 \le C\sqrt{\rho_q\rho_r}\,
 \frac{\gcd(q,r)^2}{qr}.
\]
For the canonical BEMOC populations, a finite union of affine or
quasi-affine sequences, the arithmetic estimate
\[
 \sum_{q,r\le CM}\frac{\gcd(q,r)^2}{\sqrt{qr}}=O(M^2)
\]
implies a total $O(M)=O(\sqrt N)$ error.  This is a credible repair, but it is
not the argument presently written in \texttt{sumdist.tex}.

\section{Necklace term}

\subsection{Residual estimate}

\fatal{} The proof of $A_R\lesssim1$ in lines 721--729 does not perform a
valid double summation.  A bound with no adequate decay in
$\ell=|j-k|$ is summed as if such decay were present.  The relevant vertical
separation is
\[
 |z_j-z_k|\asymp\frac{|j^2-k^2|}{N}
 =\frac{|j-k|(j+k)}N,
\]
not $j\ell/N$ uniformly.

Moreover, subtracting a single constant does not produce a mixed-second-
derivative estimate.  One needs the four-term bilinear cancellation coming
from zero mass in each variable.  Diagonal residual blocks require a separate
scaled calculation because the mixed derivative is logarithmically singular.
The assertion that cross-hemisphere pairs have separation bounded below is
false near the equator.  Therefore $A_R\lesssim1$ is unproved.

\subsection{Far Simpson pairs}

\fatal{} In lines 752--766, substitution into the displayed estimate for
$k\gg j$ gives
\[
 j^2k^{-7}N^{3/2},
\]
not $j^2k^{-7}N^{-1/2}$.  This is an $N^2$ scaling discrepancy.  A different
nonlocal kernel estimate may control this regime, but the printed calculation
does not.

The statement that cross-hemisphere interactions are ``exponentially small in
the number of vanishing moments'' is also invalid: there are four fixed
moments, and the corresponding Taylor remainder is algebraic.  Consequently
$A_S\lesssim\sqrt N$, and hence $A\lesssim\sqrt N$, do not follow from the
current proof.

\subsection{Kernel bounds and poles}

\major{} The reduced-kernel lemma correctly identifies the local singular
coefficient
\[
 \frac1{2\pi\rho^3}(z-z')^2\log|z-z'|.
\]
It does not, however, establish the uniform latitude-dependent derivative
bounds later used for the analytic coefficient and remainder.  The asserted
$\rho^{-9}$ bound for eight derivatives of the analytic part is likewise not
derived.  Such uniformity is essential near the polar bands.

\subsection{Unsupported necklace constant}

\fatal{} Lines 815--838 replace a proof by numerical assertions:
\[
 C_1\approx-7.1\times10^{-4},\qquad
 |C_\ell|\le |C_1|,
 \qquad
 \sum_{\ell\ge1}|C_\ell|\le7.2\times10^{-4}.
\]
No formulas or analytic bounds establish these inequalities.  The local
fixed-shell asymptotic is not shown to be uniform in shell number, latitude,
or the polar and equatorial transition regions.  Therefore the manuscript
does not prove:
\begin{itemize}
\item $A=\Theta(\sqrt N)$;
\item convergence of $A/\sqrt N$;
\item $c_A\approx0.0404$ or its stated interval;
\item that convergence of $D_N/\sqrt N$ reduces only to that of $C/\sqrt N$.
\end{itemize}
A positive diagonal sum alone does not lower-bound the full norm: negative
off-diagonal inner products can cancel it.

\section{Other corrections}

\begin{itemize}
\item The manuscript declares the necklace section twice consecutively, with
the same label both times.
\item The monotonicity sign in the proof of the within-ring envelope near
line 483 is reversed, although the final envelope can remain correct.
\item Calling the bands ``equal mass'' is misleading: their masses are
$r_j/N$ and vary with $j$.
\item The phrase ``main term'' for $B$ conflicts with the later claim that
$A,B,C$ all have the same order.
\item Several claims described as theorems are later said to be merely
``confirmed numerically.''  Their status must be made consistent.
\end{itemize}

\section{Logical status and repair plan}

The manuscript currently proves
\[
 D_N\ge c\sqrt N,
 \qquad
 B=\frac{2\pi(2\sqrt2-1)}9\sqrt N+o(\sqrt N),
\]
together with the exact decomposition $D_N=A+B+C$ and the signs
$A\ge0$, $B+C\ge0$.

It does not currently prove
\[
 A\lesssim\sqrt N,qquad |C|\lesssim\sqrt N,
\]
and hence does not prove $D_N\asymp\sqrt N$.

A viable repair has two components:
\begin{enumerate}
\item Replace the nonuniform branch-point argument by the bounded-variation
Fourier estimate and the resulting $\gcd$ sum.
\item Write the height error globally as a sum of signed band measures with
zero mass and zero first moment.  In the double energy, these moments
annihilate the scale-dependent quadratic term in
$(s-t)^2\log|s-t|$.  A pole-compatible near/far block decomposition should
then yield a summable estimate of the form
\[
 |(\mu_j\otimes\mu_k)(F)|
 \le C\frac{\min(j,k)}{\sqrt N}(1+|j-k|)^{-2}
\]
for comparable bands, with stronger estimates elsewhere.
\end{enumerate}

All claims concerning a limiting necklace constant should remain numerical
conjectures until the shell estimates and exchange of limits are proved
uniformly.

\begin{thebibliography}{9}
\bibitem{BEMOC}
C.~Beltr\'an, U.~Etayo, J.~Marzo and J.~Ortega-Cerd\`a,
\emph{A sequence of polynomials with optimal condition number},
J. Amer. Math. Soc. \textbf{34} (2021), 219--244.
\bibitem{Beck}
J.~Beck, \emph{Sums of distances between points on a sphere---an application
of the theory of irregularities of distribution to discrete geometry},
Mathematika \textbf{31} (1984), 33--41.
\bibitem{Stolarsky}
K.~B.~Stolarsky, \emph{Sums of distances between points on a sphere. II},
Proc. Amer. Math. Soc. \textbf{41} (1973), 575--582.
\end{thebibliography}

\end{document}
